% =====================================================================
%  1. Introdução
% =====================================================================
\section{Introdução}
\label{sec:introducao}

A avaliação automática de uma frase isolada, sem gabarito e sem referência, constitui
problema de ordem prática antes de teórica. Um sistema que redige \emph{Há três dias o
incêndio atingiu uma área de mata nativa} a partir dos campos de um registro de eventos de
fogo precisa verificar o que produziu, e não dispõe de termo de comparação: não há tradução
de origem nem resumo humano. Métricas que pressupõem um texto de referência, como BLEU e
ROUGE, não se aplicam, e a verificação recai sobre a frase em si.

Os métodos disponíveis para essa tarefa são numerosos e heterogêneos: léxicos associados a
distância de edição~\cite{damerau1964,levenshtein1966,kukich1992}, gramáticas de padrões
escritos manualmente~\cite{naber2003,milkowski2010}, analisadores
sintáticos~\cite{nivre2016,rademaker2017}, modelos de linguagem dos quais derivam medidas de
aceitabilidade como o SLOR~\cite{pauls2012,lau2017} e a
pseudo-log-verossimilhança~\cite{wang2019,salazar2020}, e vetores
distribucionais~\cite{mikolov2013}. Não respondem, contudo, à mesma pergunta: alguns
produzem veredito exato, outros veredito estimado, e outros apenas um valor numérico.

Tendo em vista um sistema que gera textos sobre eventos de fogo em linguagem natural a
partir apenas de informações numéricas extraídas de um banco de dados, é necessária uma
metodologia parametrizada e rastreável para avaliar os textos gerados. Trata-se de um
problema de múltiplas dimensões --- gramática, estrutura, ortografia e semântica.
em cada uma dessas dimensões.
